\input{../tex-inputs/latex-leseansicht-vorspann}

\section[Rainer Maria Rilke an Arthur Schnitzler, 14. 1. 1902]{L04542 Rainer Maria Rilke an Arthur Schnitzler, 14. 1. 1902}
\nopagebreak\mylabel{L04542v}
\rehead{ }\normalsize\beginnumbering\briefempfaengerindex{Schnitzler, Arthur@\textsc{Schnitzler, Arthur}!zzzRilke, Rainer Maria@\emph{von Rainer Maria Rilke}!1902-01-142@{14. 1. 1902}|(be}
\toendnotes[C]{\smallbreak\pagebreak[2]}
\correspDesc{Versand  durch Rainer Maria Rilke am 14. 1. 1902 in ?? [Bauernhaus von Clara und Rainer Maria Rilke in Westerwede]
\newline{}Erhalt  durch Arthur Schnitzler im Zeitraum [15. 1. 1902 –
                  17. 1. 1902?] in Frankgasse 1}\toendnotes[C]{\smallbreak}
\Standort{DLA, A:Schnitzler, 96.60.59.}
\physDesc{Brief, 1 Blatt, 4 Seiten, 3617 Zeichen
\newline{}Handschrift: schwarze Tinte, deutsche Kurrent
\newline{}Schnitzler: mit rotem Buntstift eine Unterstreichung 
\newline{}Ordnung: 1) Die Originale zweier Briefe und einer Karte von Rilke\pwindex{Rilke, Rainer Maria 4.\,12.\,1875 Prag – 29.\,12.\,1926 Montreux@\textsc{Rilke, Rainer Maria} (4.\,12.\,1875 Prag – 29.\,12.\,1926 Montreux)|pw} an Schnitzler fehlen in der Überlieferung im
                                 Nachlass Schnitzlers. Die
                                 Karte ist diese: XXXX Auszeichnungsfehler: Dokument L04541 nicht gefunden.
                                 Der vorliegende Briefe wird von Heinrich Schnitzler\pwindex{Schnitzler, Heinrich 9.\,8.\,1902 Hinterbrühl – 12.\,7.\,1982 Wien@\textsc{Schnitzler, Heinrich} (9.\,8.\,1902 Hinterbrühl – 12.\,7.\,1982 Wien)|pw} in der ersten Edition (S. 293) als in
                                 seinem Privatbesitz befindlich ausgewiesen.  2) Im März 1996 wurde der Brief vom Auktionshaus \emph{Stargardt} verkauft (Nr. 663, S. 117, Nr. 310) und vom \emph{Deutschen Literaturarchiv Marbach} erstanden.}
\buchAbdrucke{\weitereDrucke{\emph{Rainer Maria Rilke und Arthur Schnitzler. Ihr Briefwechsel.}. Mit Anmerkungen herausgegeben von Heinrich Schnitzler In: \emph{Wort und Wahrheit}, Jg. 13, Nr. 4, April 1958, S. 282–298, hier S. 290–291.} }\toendnotes[C]{\smallbreak}
\pstart
           \raggedleft{}{\pb}Weſterwede bei \textsc{Worpswede}{\\}(über \textsc{Bremen})\oindex{Westerwede@\textbf{Westerwede}|pw} am 14. Jan. 1902.\pend
           
\pstart{}Mein verehrter \textsc{Arthur Schnitzler},\pend\vspace{0.5em}
\pstart
           es muſs ein{ }ſchöner Abend\eventindex{Deutsches Theater Berlin@\textbf{Deutsches Theater Berlin}!Uraufführung von Lebendige Stunden, 4.1.1902@Uraufführung von Lebendige Stunden, 4.1.1902|pwv}
               geweſen{ }ſein in \textsc{Berlin\oindex{Berlin@\textbf{Berlin}|pw}}! Soweit man das aus dem
               Zeitungsdeutſch erkennen konnte, hab’ ich es erkannt, und aus einer geſchickteren
               Beſprechung kann ich{ }ſogar die Art der »letzten
                  Masken\pwindex{Schnitzler, Arthur 15. 5. 1862 Wien – 21. 10. 1931 ebd.@\textsc{Schnitzler, Arthur} (15. 5. 1862 Wien – 21. 10. 1931 ebd.)!letzten Masken@\strich\emph{Die letzten Masken}|pw}« vermuthen!\pend
           
\pstart
           Indeſſen: ich habe heute von etwas Anderem zu{ }ſprechen, von etwas Schwerem, das mich
               betroffen hat, nicht ganz unerwartet, aber immerhin mit jener Gewaltſamkeit, die jede
               Wirklichkeit vor dem Gefürchteten voraus hat.\pend
           
\pstart
           Sie wiſſen, daſs ich mir erſt ganz kürzlich ein Heim\oindex{?? [Bauernhaus von Clara und Rainer Maria Rilke in Westerwede]@\textbf{?? [Bauernhaus von Clara und Rainer Maria Rilke in Westerwede]}|pwv} in dieſer Einſamkeit geſchaffen habe und daſs meine
               Hoffnung war, in der Stille dieſes Heims\oindex{?? [Bauernhaus von Clara und Rainer Maria Rilke in Westerwede]@\textbf{?? [Bauernhaus von Clara und Rainer Maria Rilke in Westerwede]}|pwv}, einige Fortſchritte zu machen, die ich jetzt kommen
               fühlte und die mich in meiner neuen, geründeten und durch ihre Geſchloſſenheit
               wahrhaften Welt {\pb}bereiter finden{ }ſollten und fähiger mit
               ihnen zu gehen. Nun will es ein Verhängnis, daſs ich gerade jetzt einen Zuſchuſs
               verliere, von dem wir faſt ausſchließlich gelebt haben, (denn meine Bücher und
               Arbeiten tragen, trotz jahrelanger Bemühungen,\strikeout{)}{ }ſo gut wie nichts –) –{ }ſo daſs ich, da
               ein Verdienen von hier aus, das nur durch die vertrauensvolle Hülfe eines Verlegers
               möglich geweſen wäre, – ausgeſchloſſen iſt, alles verlaſſen und in die Fremde gehen
               muſs, mir und {[}den{]} meinen Brot zu holen. Denken Sie meine
               Beſtürzung: der Zuſammenhang mit der Zeit fehlt mir ganz, den faſt jede
               Beſchäftigung, die Geld einbringt, verlangt und ich weiß, daſs es für meine Kunſt
               keine größere Feindſchaft giebt als die Zeit, das Tägliche {\dots} das Allzutägliche! Trotzdem muſs etwas gefunden werden, bald, beſſer heute, als
               morgen, und ich kann froh{ }ſein, wenn es etwas iſt, was mir ermöglicht meinen Weg
               nicht ganz aufzugeben, ihn wenigstens mit den Augen feſtzuhalten, wenn {\pb}auch die Hände eine Weile von ihm laſſen müſſen. Ich
               habe mich in dieſen Tagen der Angſt nach mehreren Seiten gewandt, und es iſt nicht
               unmöglich, daſs ich nach \textsc{Wien\oindex{Wien@\textbf{Wien}|pw}} komme, um dort eine
                  \textsc{Korrespondenz} für eine Zeitſchrift zu übernehmen. Dieſer
               Poſten würde aber nicht genügen und ich müſste nebenbei noch andere regelmäßige \textsc{Korrespondenzen} und Mitarbeiterſchaften ev. auch an
               Zeitungen im Ausland oder im übrigen \textsc{Oesterreich\oindex{Österreich@\textbf{Österreich}|pw}} übernehmen dürfen (Gebiet: bildende Künſte, Litteratur, ev. auch
                  \textsc{Theater}{\dots}) wenn mir damit wirklich gedient{ }ſein{ }ſoll.
               Ich{ }ſchreibe Ihnen jedenfalls wie{ }ſich die Sache mit dem \substVorne{}\textsuperscript{\textsc{C}}\substDazwischen{}\textsc{K}\substHinten{}o\textsc{rrespondentenposten} entwickelt und bitte Sie, wenn
               Sie gelegentlich (was doch leicht möglich iſt) von \textsc{Vakanzen}
               hören,{ }ſich meiner zu erinnern: Sie thun mir einen großen Dienſt, verehrter Herr \textsc{Schnitzler}. Ich{ }ſtehe mit meinen Lieben\pwindex{Westhoff, Clara 21.\,11.\,1878 Bremen – 9.\,3.\,1954 Fischerhude@\textsc{Westhoff, Clara} (21.\,11.\,1878 Bremen – 9.\,3.\,1954 Fischerhude)|pwv}\pwindex{Sieber, Ruth 12.\,12.\,1901 Westerwede – 27.\,11.\,1972@\textsc{Sieber, Ruth} (12.\,12.\,1901 Westerwede – 27.\,11.\,1972)|pwv} an dieſem böſen
               Kreuzweg im argen Winde und warte wohin die Nothwendigkeit {\pb}mich drängt, herzlich bereit, aus dem Leben, das{ }ſie mir
               zumuthet, \textsc{mein} Leben zu machen.\pend
           
\pstart
           Meine Lieben\pwindex{Westhoff, Clara 21.\,11.\,1878 Bremen – 9.\,3.\,1954 Fischerhude@\textsc{Westhoff, Clara} (21.\,11.\,1878 Bremen – 9.\,3.\,1954 Fischerhude)|pwv}\pwindex{Sieber, Ruth 12.\,12.\,1901 Westerwede – 27.\,11.\,1972@\textsc{Sieber, Ruth} (12.\,12.\,1901 Westerwede – 27.\,11.\,1972)|pwv},
               (d. h. meine Frau\pwindex{Westhoff, Clara 21.\,11.\,1878 Bremen – 9.\,3.\,1954 Fischerhude@\textsc{Westhoff, Clara} (21.\,11.\,1878 Bremen – 9.\,3.\,1954 Fischerhude)|pwv} und
               unſere kleine Tochter\pwindex{Sieber, Ruth 12.\,12.\,1901 Westerwede – 27.\,11.\,1972@\textsc{Sieber, Ruth} (12.\,12.\,1901 Westerwede – 27.\,11.\,1972)|pwv})
               bleiben im \textsc{Westerweder} Hauſe\oindex{?? [Bauernhaus von Clara und Rainer Maria Rilke in Westerwede]@\textbf{?? [Bauernhaus von Clara und Rainer Maria Rilke in Westerwede]}|pw} zurück; trotzdem fiele die
               Nothwendigkeit eines getrennten Hausſtandes pekuniär nicht arg ins Gewicht, weil das
               Leben hier ungewöhnlich billig iſt und beide Haushalte lange nicht{ }ſoviel koſten
               würden, wie ein gemeinſames Wohnen in dem theueren Wien\oindex{Wien@\textbf{Wien}|pw}. Auch hängt die \textsc{Kunst} meiner Frau\pwindex{Westhoff, Clara 21.\,11.\,1878 Bremen – 9.\,3.\,1954 Fischerhude@\textsc{Westhoff, Clara} (21.\,11.\,1878 Bremen – 9.\,3.\,1954 Fischerhude)|pwv} zu{ }ſehr mit dieſem Lande zuſammen, als daſs ich ihr
               jetzt eine Trennung davon zumuthen dürfte. Es bleibt alſo nur ein einziger,{ }ſchmerzhafter Ausweg. Nun hoffe ich, daſs ich den Poſten in \textsc{Wien\oindex{Wien@\textbf{Wien}|pw}} bekomme, denn ich würde gern in \textsc{Wien\oindex{Wien@\textbf{Wien}|pw}}{ }ſein und glaube wohl, daſs ich mir da auch Raum{ }ſchaffen
               und Gutes leiſten könnte! Sie hören bald mehr und inzwiſchen wiſſen Sie mir
               vielleicht auch einen beſtimmten Rath zu geben{\dotsfour} daran zu
               denken, daſs ich Ihnen nahe wäre, hat viel Schönes und Liebes für mich!\pend
           \pstart Ihr: \spacefill\mbox{Rainer Maria Rilke}\pend{}\selectlanguage{ngerman}\endnumbering\briefempfaengerindex{Schnitzler, Arthur@\textsc{Schnitzler, Arthur}!zzzRilke, Rainer Maria@\emph{von Rainer Maria Rilke}!1902-01-142@{14. 1. 1902}|)be}\mylabel{L04542h}
\begin{anhang}
\end{anhang}\newcommand{\dateiname}{L04542}\newcommand{\titel}{Rainer Maria Rilke an Arthur Schnitzler, 14. 1. 1902}\newcommand{\editorInnen}{Herausgegeben von Jahnke, SelmaMüller, Martin Anton}\input{../tex-inputs/latex-leseansicht-abspann}
